% GENERATED FILE. Edit canonical lessons, then run npm run generate:books.
% canonical-source-hash: fnv1a64:75f5532fbc55c0a2
% canonical-lessons: BN-C460-nishthur, BN-C460-kothor, BN-C460-ugro, BN-C460-svarthopor, BN-C460-otyachar-kora

\chapter{Harsh Natures}
\label{ch:bn-harsh-natures}
\hlchaptermodality{460}

\begin{chapteropening}
\textbf{By the end of this chapter:} \emph{I can say that someone is cruel, stern or selfish.}

Complete the last lesson of chapter 460: I can say that someone is cruel, stern or selfish.
\end{chapteropening}

\section[\texorpdfstring{\bn{নিষ্ঠুর} (nishṭhur)}{নিষ্ঠুর (nishṭhur)}]{\bn{নিষ্ঠুর} (nishṭhur) — cruel}

\label{lesson:BN-C460-nishthur}



\begin{quote}
Before the new one: say the Bengali for ill-mannered, then the Bengali for to flatter.
\end{quote}

\subsection*{You'll want to know: \bn{নিষ্ঠুর}}
\textbf{\bn{নিষ্ঠুর}} — \emph{nishṭhur} — \textquotedblleft{}cruel\textquotedblright{}.

\begin{quote}
\bn{নিষ্ঠুর} \bn{লোকটা} \bn{কুকুরটাকে} \bn{মারল।}
\end{quote}

\begin{quote}
\emph{nishṭhur lokṭā kukurṭāke mārlô.}
\end{quote}

\begin{quote}
The cruel man beat the dog.
\end{quote}

\begin{cousinweb}
\textbf{\bn{নিষ্ঠুর}} (\emph{nishṭhur}) is a tatsama, Sanskrit \emph{niṣṭhura}, \textquotedblleft{}hard, rough, harsh, cruel\textquotedblright{} (Monier-Williams).
\end{cousinweb}

\begin{grammarlens}[title={What you've built}]
One more word for describing what you lived and felt.
\end{grammarlens}

\subsection*{Your turn}
\begin{itemize}
\raggedright
  \item \emph{Say it:} \emph{nishṭhur}
  \item \emph{Say it:} \emph{nishṭhur}, once more
  \item \emph{Say it:} \emph{nishṭhur lokṭā kukurṭāke mārlô.} — The cruel man beat the dog.
  \item \emph{Recall:} say the Bengali for ill-mannered, then the Bengali for to flatter, then say \emph{nishṭhur} again
\end{itemize}

\begin{tcolorbox}[breakable,colback=teal!4,colframe=teal!35!black,title={Before you move on}]
What does \bn{নিষ্ঠুর} mean? (\textquotedblleft{}Cruel\textquotedblright{}.) Say it once more.
\end{tcolorbox}
\section[\texorpdfstring{\bn{কঠোর} (kôṭhor)}{কঠোর (kôṭhor)}]{\bn{কঠোর} (kôṭhor) — stern; harsh}

\label{lesson:BN-C460-kothor}



\begin{quote}
Before the new one: say the Bengali for to flatter, then the Bengali for cruel.
\end{quote}

\subsection*{You'll want to know: \bn{কঠোর}}
\textbf{\bn{কঠোর}} — \emph{kôṭhor} — \textquotedblleft{}stern; harsh\textquotedblright{}.

\begin{quote}
\bn{নতুন} \bn{বিচারক} \bn{কঠোর} \bn{শাস্তি} \bn{দিলেন।}
\end{quote}

\begin{quote}
\emph{nôtun bichārôk kôṭhor shāsti dilen.}
\end{quote}

\begin{quote}
The new judge handed down a harsh punishment.
\end{quote}

\begin{cousinweb}
\textbf{\bn{কঠোর}} (\emph{kôṭhor}) is a tatsama, Sanskrit \emph{kaṭhora}, \textquotedblleft{}hard, solid, stiff\textquotedblright{} (Monier-Williams); it shows the new letter \bn{ঠ}.
\end{cousinweb}

\begin{grammarlens}[title={What you've built}]
One more word for describing what you lived and felt.
\end{grammarlens}

\subsection*{Your turn}
\begin{itemize}
\raggedright
  \item \emph{Say it:} \emph{kôṭhor}
  \item \emph{Say it:} \emph{kôṭhor}, once more
  \item \emph{Say it:} \emph{nôtun bichārôk kôṭhor shāsti dilen.} — The new judge handed down a harsh punishment.
  \item \emph{Recall:} say the Bengali for to flatter, then the Bengali for cruel, then say \emph{kôṭhor} again
\end{itemize}

\begin{tcolorbox}[breakable,colback=teal!4,colframe=teal!35!black,title={Before you move on}]
What does \bn{কঠোর} mean? (\textquotedblleft{}Stern; harsh\textquotedblright{}.) Say it once more.
\end{tcolorbox}
\section[\texorpdfstring{\bn{উগ্র} (ugrô)}{উগ্র (ugrô)}]{\bn{উগ্র} (ugrô) — aggressive; fierce}

\label{lesson:BN-C460-ugro}



\begin{quote}
Before the new one: say the Bengali for cruel, then the Bengali for stern.
\end{quote}

\subsection*{You'll want to know: \bn{উগ্র}}
\textbf{\bn{উগ্র}} — \emph{ugrô} — \textquotedblleft{}aggressive; fierce\textquotedblright{}.

\begin{quote}
\bn{উগ্র} \bn{লোকটা} \bn{রাস্তায়} \bn{ঝগড়া} \bn{শুরু} \bn{করল।}
\end{quote}

\begin{quote}
\emph{ugrô lokṭā rāstāy jhôgṛā shuru kôrlô.}
\end{quote}

\begin{quote}
The aggressive man started a quarrel in the street.
\end{quote}

\begin{cousinweb}
\textbf{\bn{উগ্র}} (\emph{ugrô}) is a tatsama, Sanskrit \emph{ugra}, \textquotedblleft{}powerful, violent, fierce\textquotedblright{} (Monier-Williams).
\end{cousinweb}

\begin{grammarlens}[title={What you've built}]
One more word for describing what you lived and felt.
\end{grammarlens}

\subsection*{Your turn}
\begin{itemize}
\raggedright
  \item \emph{Say it:} \emph{ugrô}
  \item \emph{Say it:} \emph{ugrô}, once more
  \item \emph{Say it:} \emph{ugrô lokṭā rāstāy jhôgṛā shuru kôrlô.} — The aggressive man started a quarrel in the street.
  \item \emph{Recall:} say the Bengali for cruel, then the Bengali for stern, then say \emph{ugrô} again
\end{itemize}

\begin{tcolorbox}[breakable,colback=teal!4,colframe=teal!35!black,title={Before you move on}]
What does \bn{উগ্র} mean? (\textquotedblleft{}Aggressive; fierce\textquotedblright{}.) Say it once more.
\end{tcolorbox}
\section[\texorpdfstring{\bn{স্বার্থপর} (svārthôpôr)}{স্বার্থপর (svārthôpôr)}]{\bn{স্বার্থপর} (svārthôpôr) — selfish}

\label{lesson:BN-C460-svarthopor}



\begin{quote}
Before the new one: say the Bengali for stern, then the Bengali for aggressive.
\end{quote}

\subsection*{You'll want to know: \bn{স্বার্থপর}}
\textbf{\bn{স্বার্থপর}} — \emph{svārthôpôr} — \textquotedblleft{}selfish\textquotedblright{}.

\begin{quote}
\bn{স্বার্থপর} \bn{ভাইটা} \bn{একা} \bn{পুরো} \bn{কেকটা} \bn{খেল।}
\end{quote}

\begin{quote}
\emph{svārthôpôr bhāiṭā ekā puro kekṭā khelô.}
\end{quote}

\begin{quote}
The selfish brother ate the whole cake by himself.
\end{quote}

\begin{cousinweb}
\textbf{\bn{স্বার্থপর}} (\emph{svārthôpôr}) is a tatsama, Sanskrit \emph{svārthapara}, \textquotedblleft{}intent on one's own advantage\textquotedblright{} (Monier-Williams): \emph{sva} \textquotedblleft{}own\textquotedblright{}, \emph{artha} \textquotedblleft{}gain\textquotedblright{}, \emph{para} \textquotedblleft{}devoted to\textquotedblright{}.
\end{cousinweb}

\begin{grammarlens}[title={What you've built}]
One more word for describing what you lived and felt.
\end{grammarlens}

\subsection*{Your turn}
\begin{itemize}
\raggedright
  \item \emph{Say it:} \emph{svārthôpôr}
  \item \emph{Say it:} \emph{svārthôpôr}, once more
  \item \emph{Say it:} \emph{svārthôpôr bhāiṭā ekā puro kekṭā khelô.} — The selfish brother ate the whole cake by himself.
  \item \emph{Recall:} say the Bengali for stern, then the Bengali for aggressive, then say \emph{svārthôpôr} again
\end{itemize}

\begin{tcolorbox}[breakable,colback=teal!4,colframe=teal!35!black,title={Before you move on}]
What does \bn{স্বার্থপর} mean? (\textquotedblleft{}Selfish\textquotedblright{}.) Say it once more.
\end{tcolorbox}
\section[\texorpdfstring{\bn{অত্যাচার} \bn{করা} (ôtyāchār kôrā)}{অত্যাচার করা (ôtyāchār kôrā)}]{\bn{অত্যাচার} \bn{করা} (ôtyāchār kôrā) — to oppress; to torment}

\label{lesson:BN-C460-otyachar-kora}



\begin{quote}
Before the new one: say the Bengali for aggressive, then the Bengali for selfish.
\end{quote}

\subsection*{You'll want to know: \bn{অত্যাচার} \bn{করা}}
\textbf{\bn{অত্যাচার} \bn{করা}} — \emph{ôtyāchār kôrā} — \textquotedblleft{}to oppress; to torment\textquotedblright{}.

\begin{quote}
\bn{নিষ্ঠুর} \bn{রাজা} \bn{গরিবদের} \bn{ওপর} \bn{অত্যাচার} \bn{করতেন।}
\end{quote}

\begin{quote}
\emph{nishṭhur rājā gôribder opôr ôtyāchār kôrten.}
\end{quote}

\begin{quote}
The cruel king used to oppress the poor.
\end{quote}

Past (he/she …): \bn{তিনি} \bn{গ্রামের} \bn{লোকদের} \bn{ওপর} \bn{অত্যাচার} \bn{করলেন}, \emph{tini grāmer lokder opôr ôtyāchār kôrlen}.

\begin{cousinweb}
\textbf{\bn{অত্যাচার} \bn{করা}} (\emph{ôtyāchār kôrā}) is Sanskrit \emph{atyācāra}, which Monier-Williams gives for going beyond the usual practice; Bengali uses it for oppression, a sense the dictionary does not give. \emph{kôrā} is \textquotedblleft{}to do\textquotedblright{} (CDIAL 2814).
\end{cousinweb}

\begin{grammarlens}[title={What you've built}]
One more word for describing what you lived and felt.
\end{grammarlens}

\subsection*{Your turn}
\begin{itemize}
\raggedright
  \item \emph{Say it:} \emph{ôtyāchār kôrā}
  \item \emph{Say it:} \emph{ôtyāchār kôrā}, once more
  \item \emph{Say it:} \emph{nishṭhur rājā gôribder opôr ôtyāchār kôrten.} — The cruel king used to oppress the poor.
  \item \emph{Recall:} say the Bengali for aggressive, then the Bengali for selfish, then say \emph{ôtyāchār kôrā} again
\end{itemize}

\begin{tcolorbox}[breakable,colback=teal!4,colframe=teal!35!black,title={Before you move on}]
What does \bn{অত্যাচার} \bn{করা} mean? (\textquotedblleft{}To oppress; to torment\textquotedblright{}.) Say it once more.
\end{tcolorbox}
